\section{Physical cycle obstructions for an actual observation map}
\label{sec:physical-cycle}
This part closes a limitation of the scalar four-cycle example: the
obstruction is now realized inside the actual divergence-zero flat physical
cochain domain. The kernel estimate below is proved for all actual cochains,
not just the special fields used as negative witnesses. The conclusion
concerns this observation and this fixed penalty. It does not obstruct all
partitions or the optimal-rank question from the preceding parts.

\subsection{The actual form and rectangular observations}
Fix $D\geq3$, $N_c\geq2$ and $X=(\mathbb Z/4\mathbb Z)^D$.
Write $V=\mathbb R^{N_c^2-1}$ with its standard Euclidean norm and
$A=(A_\ell(x))_{x\in X,1\leq\ell\leq D}$ for an actual one-cochain.
All norms use counting measure. Choose the actual total coefficients
\[
 \alpha=(1,6,4,1),\qquad a(x)=\alpha_{x_1},\qquad
 w(i,x,\ell)=\alpha_{x_1}-1
\]
for every derivative direction $i$ and physical component $\ell$.
Define actual forward and backward differences by
\[
 D_i^+f(x)=f(x+e_i)-f(x),\qquad D_i^-f(x)=f(x)-f(x-e_i).
\]
The flat curl and backward divergence, with one curl for each unordered
pair, are
\[
 (\operatorname{curl}A)_{ij}=D_i^+A_j-D_j^+A_i\quad(i<j),\qquad
 \operatorname{div}^-A=\sum_iD_i^-A_i.
\]
Their counting quadratic form and the original componentwise repair are
\begin{align*}
 H_0(A)&=\sum_{i<j,x}\|D_i^+A_j(x)-D_j^+A_i(x)\|^2
       +\sum_x\Big\|\sum_iD_i^-A_i(x)\Big\|^2,\\
 R_w(A)&=\sum_{i,x,\ell}w(i,x,\ell)\|D_i^+A_\ell(x)\|^2.
\end{align*}
In operator notation the repair is $\langle A,G^*WG A\rangle$ for the
actual stacked gradient $G$ and diagonal multiplication $W$.
These formulas fix the energy on paper; no comparison certificate is supplied.

\paragraph{Deriving the flat identity.}
Periodic reindexing gives
\[
 \sum_x\langle D_i^+f,h\rangle=-\sum_x\langle f,D_i^-h\rangle.
\]
Coordinate shifts commute. Applying this formula twice and real symmetry
therefore yields
\[
 \sum_x\langle D_i^+f,D_j^+g\rangle
 =\sum_x\langle D_j^-f,D_i^-g\rangle.
\]
Also $\sum_x\|D_i^-f\|^2=\sum_x\|D_i^+f\|^2$ by translation.
Expand each curl square with $f=A_j$, $g=A_i$: it supplies the two
off-diagonal gradient norms and the negative mixed terms. Expansion of
the divergence square supplies the diagonal norms and the opposite
mixed terms. Cancellation for each $i<j$ proves, for every cochain,
\[
 H_0(A)=\sum_{i,x,\ell}\|D_i^+A_\ell(x)\|^2.
\]
No divergence condition or lift ansatz was used. Define $H=H_0+R_w$.
Adding the displayed repair, with $1+w(i,x,\ell)=\alpha_{x_1}$, gives
\begin{equation}\label{eq:actual-cycle-energy}
 H(A)=\sum_{x\in X}\sum_{i=1}^D\sum_{\ell=1}^D
 \alpha_{x_1}\|A_\ell(x+e_i)-A_\ell(x)\|^2.
\end{equation}
The cancellation was unweighted; the varying repair is added afterwards.
We do not move $\alpha$ through summation by parts or assert a weighted
curl--divergence identity. This is the paper counterpart of
\texttt{finite\_flat\_curl\_div\_energy} and
\texttt{componentwiseWeightedRepair\_energy} in the pinned formal source.
Thus the coefficient is the original shared scalar total coefficient,
not a substituted energy or an interacting Hessian.

Put $I_0=\{0,1\}$ and $I_1=\{2,3\}$. For every transverse
$y\in(\mathbb Z/4\mathbb Z)^{D-1}$ the blocks $I_b\times\{y\}$
are nonempty axis-aligned rectangles. They are disjoint and cover $X$.
For each physical component $\ell$ and each standard internal coordinate $k$,
define the actual normalized means
\begin{equation}\label{eq:actual-cycle-C}
 (CA)_{b,y,\ell,k}=\frac{1}{\sqrt2}\sum_{t\in I_b} A_\ell(t,y)_k.
\end{equation}
The output includes \emph{all} these coordinates. If
$h=D(N_c^2-1)$, its dimension is $2h4^{D-1}$; no optimality is asserted.

The spatial mean of $a$ is $\mu=3$, and $1\leq a\leq2\mu$.
For $p=D/2$ its original normalized inverse moment is
\begin{equation}\label{eq:actual-cycle-moment}
 \frac1{|X|}\sum_{x\in X}\left(\frac{\mu}{a(x)}\right)^p
   =\frac{2\cdot3^p+(1/2)^p+(3/4)^p}{4}=:M_D<\infty.
\end{equation}
This is an actual member of the stated bounded-moment coefficient class
when the allowed bound is at least $M_D$. Its fixed finite size and the
large output dimension do not disprove an asymptotic optimal-rank order.

\subsection{A fixed-penalty counterexample with a proved kernel premise}
For a real cyclic vector $u=(u_0,u_1,u_2,u_3)$ set
\begin{align*}
 E(u)&=(u_1-u_0)^2+6(u_2-u_1)^2+4(u_3-u_2)^2+(u_0-u_3)^2,\\
 |u|^2&=\sum_{t=0}^3u_t^2,\qquad
 O(u)=\frac{(u_0+u_1)^2+(u_2+u_3)^2}{2}.
\end{align*}
Choose a unit internal coordinate vector $T$. The actual lift is
\[
 (Ju)_2(x)=u_{x_1}T,\qquad (Ju)_\ell(x)=0\quad(\ell\ne2).
\]
Its backward divergence is zero: its only nonzero component is constant
in coordinate $2$. It is injective and direct summation gives
\begin{equation}\label{eq:cycle-transports}
 \|Ju\|^2=4^{D-1}|u|^2,\quad H(Ju)=4^{D-1}E(u),\quad
 \|CJu\|^2=4^{D-1}O(u).
\end{equation}
The last equality uses the entire output in \eqref{eq:actual-cycle-C},
not only a selected nonzero coordinate.

\paragraph{Kernel premise on all actual fields.}
For $v=(x,-x,y,-y)$, direct algebra yields
\begin{align}
 E(v)&=11x^2+14xy+23y^2,\qquad |v|^2=2x^2+2y^2,\notag\\
 E(v)-\frac{19}{5}|v|^2
   &=\frac{17}{5}\left(x+\frac{35}{17}y\right)^2
       +\frac{84}{85}y^2\geq0.\label{eq:cycle-kernel-square}
\end{align}
If $CA=0$, both pair sums vanish for every actual scalar profile
$t\mapsto A_\ell(t,y)_k$. Apply \eqref{eq:cycle-kernel-square}
to each such profile and sum over all $y,\ell,k$. These profiles account
for the full norm of $A$. Their energies are exactly the $i=1$ term
of \eqref{eq:actual-cycle-energy}; every other derivative term is
nonnegative because $\alpha_t\geq1$. Therefore
\begin{equation}\label{eq:full-physical-kernel}
 CA=0\quad\Longrightarrow\quad H(A)\geq\frac{19}{5}\|A\|^2
 \qquad\text{for every actual cochain }A.
\end{equation}
No divergence condition, scalar-lift ansatz or supplied coercivity
certificate was used in this proof.

\paragraph{Failure on an actual divergence-zero field.}
For $u=(19,-9,-5,-5)/14$ one has
\[
 E(u)=\frac{52}{7},\qquad |u|^2=\frac{123}{49},\qquad
 O(u)=\frac{25}{49}.
\]
Consequently the actual divergence-zero field $Ju$ satisfies
\begin{equation}\label{eq:physical-negative-defect}
 H(Ju)+\frac{19}{5}\|CJu\|^2-\frac{19}{5}\|Ju\|^2
       =-\frac{6}{35}4^{D-1}<0.
\end{equation}
Equations \eqref{eq:full-physical-kernel} and
\eqref{eq:physical-negative-defect} prove the complete fixed-map statement:
kernel coercivity does not imply the same coefficient as a fixed coarse
penalty, even on the actual flat physical divergence-zero domain.
They do not rule out larger penalties or different observations.

\subsection{The exact finite penalty at the subsharp target}
The failure just proved concerns $K=19/5$. Here we determine exactly
which larger penalties work for this same observation and target.
This subsection is a paper proof with an exact rational coefficient
check; it is not part of the paired Lean theorem below.

Every $u\in\mathbb R^4$ has the unique decomposition
\[
 u=m(1,1,1,1)+(x,-x,y,-y)
       +(\delta/2,\delta/2,-\delta/2,-\delta/2),
\]
where $m=(\sum_tu_t)/4$, $x=(u_0-u_1)/2$, $y=(u_2-u_3)/2$
and $\delta=(u_0+u_1-u_2-u_3)/2$. Direct expansion gives
\begin{align*}
 |u|^2&=4m^2+2x^2+2y^2+\delta^2,\qquad O(u)=4m^2+\delta^2,\\
 E(u)&=11x^2+14xy+23y^2+7\delta^2-10\delta(x+y).
\end{align*}
Set $X=x-25\delta/2$ and $Y=y+75\delta/14$. For every real $K$,
\begin{align}
 E(u)+KO(u)-\frac{19}{5}|u|^2
 &=4\left(K-\frac{19}{5}\right)m^2
   +\frac{17}{5}\left(X+\frac{35}{17}Y\right)^2\notag\\
 &\quad+\frac{84}{85}Y^2
   +\left(K-\frac{1138}{35}\right)\delta^2.
 \label{eq:exact-subsharp-square}
\end{align}
Since $1138/35>19/5$, this proves sufficiency of $K\geq1138/35$
on every scalar profile. Necessity follows from
\[
 u_*=(13,-12,-41/7,34/7),\qquad
 O(u_*)=1,\quad |u_*|^2=\frac{18174}{49},\quad
 E(u_*)=\frac{67468}{49}.
\]
It has $m=0$, $\delta=1$ and $X=Y=0$, hence defect $K-1138/35$.

Apply the scalar inequality to every profile
$t\mapsto A_\ell(t,y)_k$ and sum over all $y,\ell,k$.
Their squared norms and observations exhaust the actual full norms,
and their energies exhaust the first-direction contribution in
\eqref{eq:actual-cycle-energy}. The remaining directions are
nonnegative. Conversely $Ju_*$ is actual divergence zero and multiplies
the scalar defect by $4^{D-1}$. Therefore the following are equivalent:
\begin{enumerate}
 \item $K\geq1138/35$;
 \item $H(A)+K\|CA\|^2\geq(19/5)\|A\|^2$ for every actual cochain;
 \item the same inequality for every actual divergence-zero cochain.
\end{enumerate}
This is a necessity-and-sufficiency statement for the fixed map, not
an obstruction to different partitions. In particular, a larger finite
penalty repairs the subsharp target even though $K=19/5$ fails.
The sharp kernel target in the next subsection behaves differently.

A second calculation records the classical block elimination behind the
constant. The kernel residual matrix and mixed vector are
\[
 M=\begin{pmatrix}17/5&7\\7&77/5\end{pmatrix},\quad
 \det M=84/25,\quad b=(5,5),\quad
 M^{-1}b=(25/2,-75/14).
\]
Thus $b^TM^{-1}b=250/7$ and the residual coefficient is
$K+16/5-250/7=K-1138/35$, agreeing with
\eqref{eq:exact-subsharp-square}. No new Schur-complement principle is claimed.

The unit151 exact verifier constructs $E$, $O$ and $|u|^2$ from the actual
four-coordinate linear forms and compares every coefficient of the
polynomial in $(m,x,y,\delta,K)$, rather than sampling values. Normal and
optimized Python in one own Colab Linux runtime each accepted five checks
and rejected three real coefficient mutations; both exits were zero and
their JSON outputs were identical. The 8794-byte archive has ten members
and nine checked inner hashes. It and the full hand derivation are
included in the reproduction package. The driver took about 0.56 seconds;
the runtime was released after recovery and exact billed consumption is
unknown. This check is neither a new Lean theorem nor independently
authored verification. Its source is
\begin{quote}\small\ttfamily
564b84230478cb8d40505cb67985470e9bae699d.
\end{quote}
\subsection{Sharp free-penalty extension and its distinct status}
The following is an analytic continuation, outside the paired Lean result.
Put $s=\sqrt{85}$ and $c_*=(17-s)/2$. Since $9<s<10$, the exact identity
\[
 E(v)-c_*|v|^2=(s-6)\left(x+\frac{7y}{s-6}\right)^2
\]
holds on the same block kernel. The nonzero vector
$v_*=(7,-7,6-s,s-6)$ attains equality. The all-profile argument above
also establishes the sharp kernel constant for the actual form in paper
proof, and the lift establishes sharpness.

For $p_\delta=(\delta/2,\delta/2,-\delta/2,-\delta/2)$, the two vectors
are orthogonal, $Cv_*=0$, and $O(p_\delta)=|p_\delta|^2=\delta^2$.
For any finite $K\geq0$, direct expansion gives
\begin{align*}
 &E(v_*+p_\delta)+K O(v_*+p_\delta)-c_*|v_*+p_\delta|^2\\
 &\hspace{18mm}=(7+K-c_*)\delta^2-10(13-s)\delta.
\end{align*}
The positive number $q_K=7+K-c_*$ permits the choice
$\delta=5(13-s)/q_K$, giving the strictly negative value
$-25(13-s)^2/q_K$. Lifting multiplies this complete defect by
$4^{D-1}$ and preserves actual divergence zero. Thus no finite quadratic
penalty repairs this same observation at the sharp kernel endpoint.
A penalty $\Phi(z)\leq K\|z\|^2$ only near zero is also defeated:
scale the negative witness by a sufficiently small nonzero real number.
Both radius conditions can be met, and its defect remains a negative
multiple of the square of that scale. This is not a strict-margin claim.

These exact analytic identities and arguments are supplied here, but their
complete physical sharp-endpoint/free-penalty statement has \emph{not}
been verified in Lean in the paired execution described next. Finite
checks in dimensions three and four, including $K=0,1,100$, are diagnostics
only, not the universal theorem or independent review.

\subsection{Formal scope and the paired reproduction record}
The fixed-penalty result, including the actual full-kernel premise and
all transport/comparison identities used to connect it, was compiled in
two separately allocated Colab Linux CPU sessions with high RAM. Each
started with a fresh clone at the same immutable source
\begin{quote}\small\ttfamily
0c706827dcbb11859e56bb6c7dc981eb37a02e69.
\end{quote}
The final theorem is named
\texttt{physicalCycleFour\_kernel\_fixed\_penalty\_counterexample}
in \texttt{PhysicalCycleFullForm.lean}. It asserts the proved
all-cochain kernel bound together with failure of the fixed-penalty
inequality on the actual divergence-zero domain. Four preceding printed
endpoints identify longitudinal energy, the full actual energy sum,
its comparison, and the kernel estimate. They do not assume transport
certificates or posit a structure containing the conclusion.

All five actual dependency readouts in each clone contain only
\texttt{propext}, \texttt{Classical.choice} and \texttt{Quot.sound}.
The normal and optimized validator in each clone accepted all five
readouts and rejected three actual oracle mutations. All actual decimal
exit sentinels were zero. Lean was version 4.29.0-rc6 with Mathlib pinned
in the included manifest. The recovered theorem source, oracle source,
compiled binary, oracle output, validator source, normal and optimized
outputs and driver have identical hashes across the two fresh executions.
The raw per-run reports still say first-clone pass and paired false:
the identical driver executes one clone per session. They remain
unmodified; the separate aggregate records the completed pair.

The first successful driver took about 374.31 seconds and the second
400.28 seconds. These are not billed connection times. Both environments
were disconnected after complete evidence recovery. No Lean compilation,
oracle or numerical campaign ran on Windows. Separate fresh executions
are not independently authored proofs or independent scientific review.
The previous failed attempt, its exact source and full compiler log are
preserved: mixed finite-index inference and nested sum type inference
were corrected without changing the theorem statements.

The package includes both public evidence archives. Their comparison is
recorded in the aggregate file:
\begin{quote}\small\ttfamily
paired-full-form148.json
\end{quote}
Full hashes and source bindings are in \texttt{inputs.json}.
Each archive has 45 checked inner manifest
members and excludes third-party installers and downloaded binary assets.
Their complete originals, including those assets, are retained in the
private evidence record. The compiled binary hash begins
\texttt{b0f8ff8654d79d02}, and the full value is in the paired comparison file.

\paragraph{Attribution and scope of novelty.}
The scalar square identities and negative vectors predate this physical
embedding in the dossier. The present consolidation adds the explicit
actual-domain embedding and the formally derived all-field comparison;
it claims no new Schur-complement principle, priority or exhaustive novelty
finding. Carlson, Haynsworth and Markham's classical paper,
\emph{A Generalization of the Schur Complement by Means of the Moore--Penrose
Inverse}, SIAM Journal on Applied Mathematics 26 (1974), 169--175,
\href{https://doi.org/10.1137/0126013}{doi:10.1137/0126013}, supplies algebraic
context. The publisher abstract and bibliography were inspected in the
source investigation; the full paywalled article was not read. The direct
finite proof here carries the mathematical claim.
